\documentclass[11pt,a4paper]{article}

\usepackage[utf8]{inputenc} % allow utf-8 input
\usepackage[T1]{fontenc}    % use 8-bit T1 fonts
\usepackage{hyperref}       % hyperlinks
\usepackage{url}            % simple URL typesetting
\usepackage{booktabs}       % professional-quality tables
\usepackage{amsfonts}       % blackboard math symbols
\usepackage{nicefrac}       % compact symbols for 1/2, etc.
\usepackage{microtype}      % microtypography
\usepackage{xcolor}         % colors
\usepackage{graphicx}       % graphics
\usepackage[normalem]{ulem} % strikethrough
\title{Duck}

\begin{document}

\maketitle

\begin{itemize}
\item [Main page](/wiki/Main\_Page)
\item [Contents](/wiki/Wikipedia:Contents)
\item [Current events](/wiki/Portal:Current\_events)
\item [Random article](/wiki/Special:Random)
\item [About Wikipedia](/wiki/Wikipedia:About)
\item [Contact us](//en.wikipedia.org/wiki/Wikipedia:Contact\_us)
\end{itemize}

\begin{itemize}
\item [Help](/wiki/Help:Contents)
\item [Learn to edit](/wiki/Help:Introduction)
\item [Community portal](/wiki/Wikipedia:Community\_portal)
\item [Recent changes](/wiki/Special:RecentChanges)
\item [Upload file](/wiki/Wikipedia:File\_upload\_wizard)
\end{itemize}

\begin{itemize}
\item [Donate](https://donate.wikimedia.org/wiki/Special:FundraiserRedirector?utm\_source=donate\&utm\_medium=sidebar\&utm\_campaign=C13\_en.wikipedia.org\&uselang=en)
\end{itemize}

\begin{itemize}
\item [Create account](/w/index.php?title=Special:CreateAccount\&returnto=Duck)
\item [Log in](/w/index.php?title=Special:UserLogin\&returnto=Duck)
\end{itemize}

\begin{itemize}
\item [Create account](/w/index.php?title=Special:CreateAccount\&returnto=Duck)
\item [Log in](/w/index.php?title=Special:UserLogin\&returnto=Duck)
\end{itemize}

\begin{itemize}
\item [Contributions](/wiki/Special:MyContributions)
\item [Talk](/wiki/Special:MyTalk)
\end{itemize}

\section{Contents}

move to sidebar

hide

\begin{itemize}
\item [(Top)](\#)
\item [1 Etymology](\#Etymology)
\item [2 Taxonomy](\#Taxonomy)
\item [3 Morphology](\#Morphology)
\item [4 Distribution and habitat](\#Distribution\_and\_habitat)
\item [5 Behaviour](\#Behaviour) Toggle Behaviour subsection
  \begin{itemize}
\item [5.1 Feeding](\#Feeding)
\item [5.2 Breeding](\#Breeding)
\item [5.3 Communication](\#Communication)
\item [5.4 Predators](\#Predators)
  \end{itemize}
\item [6 Relationship with humans](\#Relationship\_with\_humans) Toggle Relationship with humans subsection
  \begin{itemize}
\item [6.1 Hunting](\#Hunting)
\item [6.2 Domestication](\#Domestication)
\item [6.3 Heraldry](\#Heraldry)
\item [6.4 Cultural references](\#Cultural\_references)
  \end{itemize}
\item [7 See also](\#See\_also)
\item [8 Notes](\#Notes) Toggle Notes subsection
  \begin{itemize}
\item [8.1 Citations](\#Citations)
\item [8.2 Sources](\#Sources)
  \end{itemize}
\item [9 External links](\#External\_links)
\end{itemize}

- [ ] Toggle the table of contents

- [ ] 136 languages

\begin{itemize}
\item [Acèh](https://ace.wikipedia.org/wiki/It\%C3\%A9k)
\item [Afrikaans](https://af.wikipedia.org/wiki/Eend)
\item [Alemannisch](https://als.wikipedia.org/wiki/Ente)
\item [አማርኛ](https://am.wikipedia.org/wiki/\%E1\%8B\%B3\%E1\%8A\%AD\%E1\%8B\%AC)
\item [Ænglisc](https://ang.wikipedia.org/wiki/Ened)
\item [العربية](https://ar.wikipedia.org/wiki/\%D8\%A8\%D8\%B7)
\item [Aragonés](https://an.wikipedia.org/wiki/Anade)
\item [ܐܪܡܝܐ](https://arc.wikipedia.org/wiki/\%DC\%92\%DC\%9B\%DC\%90)
\item [Armãneashti](https://roa-rup.wikipedia.org/wiki/Paphi)
\item [Asturianu](https://ast.wikipedia.org/wiki/Cor\%C3\%ADu)
\item [Atikamekw](https://atj.wikipedia.org/wiki/Cicip)
\item [Авар](https://av.wikipedia.org/wiki/\%D0\%9E\%D1\%80\%D0\%B4\%D0\%B5\%D0\%BA)
\item [Aymar aru](https://ay.wikipedia.org/wiki/Unkalla)
\item [تۆرکجه](https://azb.wikipedia.org/wiki/\%D8\%A7\%D8\%A4\%D8\%B1\%D8\%AF\%DA\%A9)
\item [Basa Bali](https://ban.wikipedia.org/wiki/B\%C3\%A9b\%C3\%A9k)
\item [বাংলা](https://bn.wikipedia.org/wiki/\%E0\%A6\%B9\%E0\%A6\%BE\%E0\%A6\%81\%E0\%A6\%B8)
\item [閩南語 / Bân-lâm-gú](https://zh-min-nan.wikipedia.org/wiki/Ah)
\item [Беларуская](https://be.wikipedia.org/wiki/\%D0\%9A\%D0\%B0\%D1\%87\%D0\%BA\%D1\%96)
\item [Беларуская (тарашкевіца)](https://be-tarask.wikipedia.org/wiki/\%D0\%9A\%D0\%B0\%D1\%87\%D0\%BA\%D1\%96)
\item [Bikol Central](https://bcl.wikipedia.org/wiki/Itik)
\item [Български](https://bg.wikipedia.org/wiki/\%D0\%9F\%D0\%B0\%D1\%82\%D0\%B8\%D1\%86\%D0\%B0)
\item [Brezhoneg](https://br.wikipedia.org/wiki/Houad\_(evn))
\item [Буряад](https://bxr.wikipedia.org/wiki/\%D0\%9D\%D1\%83\%D0\%B3\%D0\%B0h\%D0\%B0\%D0\%BD)
\item [Català](https://ca.wikipedia.org/wiki/\%C3\%80necs)
\item [Чӑвашла](https://cv.wikipedia.org/wiki/\%D0\%9A\%C4\%83\%D0\%B2\%D0\%B0\%D0\%BA\%D0\%B0\%D0\%BB\%D1\%81\%D0\%B5\%D0\%BC)
\item [Čeština](https://cs.wikipedia.org/wiki/Kachna)
\item [ChiShona](https://sn.wikipedia.org/wiki/Dhadha)
\item [Cymraeg](https://cy.wikipedia.org/wiki/Hwyaden)
\item [Dagbanli](https://dag.wikipedia.org/wiki/Gbunya\%C9\%A3u)
\item [Dansk](https://da.wikipedia.org/wiki/\%C3\%86nder)
\item [Deitsch](https://pdc.wikipedia.org/wiki/Ent)
\item [Deutsch](https://de.wikipedia.org/wiki/Enten)
\item [डोटेली](https://dty.wikipedia.org/wiki/\%E0\%A4\%B9\%E0\%A4\%BE\%E0\%A4\%81\%E0\%A4\%B8)
\item [Ελληνικά](https://el.wikipedia.org/wiki/\%CE\%A0\%CE\%AC\%CF\%80\%CE\%B9\%CE\%B1)
\item [Emiliàn e rumagnòl](https://eml.wikipedia.org/wiki/An\%C3\%A0dra)
\item [Español](https://es.wikipedia.org/wiki/Pato)
\item [Esperanto](https://eo.wikipedia.org/wiki/Anaso)
\item [Euskara](https://eu.wikipedia.org/wiki/Ahate)
\item [فارسی](https://fa.wikipedia.org/wiki/\%D9\%85\%D8\%B1\%D8\%BA\%D8\%A7\%D8\%A8\%DB\%8C)
\item [Français](https://fr.wikipedia.org/wiki/Canard)
\item [Gaeilge](https://ga.wikipedia.org/wiki/Lacha)
\item [Galego](https://gl.wikipedia.org/wiki/Pato)
\item [ГӀалгӀай](https://inh.wikipedia.org/wiki/\%D0\%91\%D0\%BE\%D0\%B0\%D0\%B1\%D0\%B0\%D1\%88\%D0\%BA\%D0\%B0\%D1\%88)
\item [贛語](https://gan.wikipedia.org/wiki/\%E9\%B4\%A8)
\item [گیلکی](https://glk.wikipedia.org/wiki/\%D8\%A8\%D9\%8A\%D9\%84\%D9\%8A)
\item [𐌲𐌿𐍄𐌹𐍃𐌺](https://got.wikipedia.org/wiki/\%F0\%90\%8C\%B0\%F0\%90\%8C\%BD\%F0\%90\%8C\%B0\%F0\%90\%8C\%B8\%F0\%90\%8D\%83)
\item [गोंयची कोंकणी / Gõychi Konknni](https://gom.wikipedia.org/wiki/\%E0\%A4\%AC\%E0\%A4\%A6\%E0\%A4\%95)
\item [客家語 / Hak-kâ-ngî](https://hak.wikipedia.org/wiki/Ap-\%C3\%A8)
\item [한국어](https://ko.wikipedia.org/wiki/\%EC\%98\%A4\%EB\%A6\%AC)
\item [Hausa](https://ha.wikipedia.org/wiki/Agwagwa)
\item [Հայերեն](https://hy.wikipedia.org/wiki/\%D4\%B2\%D5\%A1\%D5\%A4\%D5\%A5\%D6\%80)
\item [हिन्दी](https://hi.wikipedia.org/wiki/\%E0\%A4\%AC\%E0\%A4\%A4\%E0\%A5\%8D\%E0\%A4\%A4\%E0\%A4\%96)
\item [Hrvatski](https://hr.wikipedia.org/wiki/Patka)
\item [Ido](https://io.wikipedia.org/wiki/Anado)
\item [Bahasa Indonesia](https://id.wikipedia.org/wiki/Itik)
\item [Iñupiatun](https://ik.wikipedia.org/wiki/Mitiq)
\item [Íslenska](https://is.wikipedia.org/wiki/\%C3\%96nd)
\item [Italiano](https://it.wikipedia.org/wiki/Anatra)
\item [עברית](https://he.wikipedia.org/wiki/\%D7\%91\%D7\%A8\%D7\%95\%D7\%95\%D7\%96)
\item [Jawa](https://jv.wikipedia.org/wiki/B\%C3\%A8b\%C3\%A8k)
\item [ಕನ್ನಡ](https://kn.wikipedia.org/wiki/\%E0\%B2\%AC\%E0\%B2\%BE\%E0\%B2\%A4\%E0\%B3\%81\%E0\%B2\%95\%E0\%B3\%8B\%E0\%B2\%B3\%E0\%B2\%BF)
\item [Kapampangan](https://pam.wikipedia.org/wiki/Bibi)
\item [ქართული](https://ka.wikipedia.org/wiki/\%E1\%83\%98\%E1\%83\%AE\%E1\%83\%95\%E1\%83\%94\%E1\%83\%91\%E1\%83\%98)
\item [कॉशुर / کٲشُر](https://ks.wikipedia.org/wiki/\%D8\%A8\%D9\%8E\%D8\%B7\%D9\%8F\%D8\%AE)
\item [Қазақша](https://kk.wikipedia.org/wiki/\%D2\%AE\%D0\%B9\%D1\%80\%D0\%B5\%D0\%BA)
\item [Ikirundi](https://rn.wikipedia.org/wiki/Imbata)
\item [Kongo](https://kg.wikipedia.org/wiki/Kivadangu)
\item [Kreyòl ayisyen](https://ht.wikipedia.org/wiki/Kanna)
\item [Кырык мары](https://mrj.wikipedia.org/wiki/\%D0\%9B\%D1\%8B\%D0\%B4\%D1\%8B\%D0\%B2\%D0\%BB\%D3\%93)
\item [ລາວ](https://lo.wikipedia.org/wiki/\%E0\%BB\%80\%E0\%BA\%9B\%E0\%BA\%B1\%E0\%BA\%94)
\item [Latina](https://la.wikipedia.org/wiki/Anas\_(avis))
\item [Latviešu](https://lv.wikipedia.org/wiki/P\%C4\%ABle)
\item [Lietuvių](https://lt.wikipedia.org/wiki/Antis)
\item [Li Niha](https://nia.wikipedia.org/wiki/Bebe)
\item [Ligure](https://lij.wikipedia.org/wiki/Annia)
\item [Limburgs](https://li.wikipedia.org/wiki/Aenj)
\item [Lingála](https://ln.wikipedia.org/wiki/Libat\%C3\%A1)
\item [Malagasy](https://mg.wikipedia.org/wiki/Ganagana)
\item [മലയാളം](https://ml.wikipedia.org/wiki/\%E0\%B4\%A4\%E0\%B4\%BE\%E0\%B4\%B1\%E0\%B4\%BE\%E0\%B4\%B5\%E0\%B5\%8D)
\item [मराठी](https://mr.wikipedia.org/wiki/\%E0\%A4\%AC\%E0\%A4\%A6\%E0\%A4\%95)
\item [مازِرونی](https://mzn.wikipedia.org/wiki/\%D8\%B3\%DB\%8C\%DA\%A9\%D8\%A7)
\item [Bahasa Melayu](https://ms.wikipedia.org/wiki/Itik)
\item [ꯃꯤꯇꯩ ꯂꯣꯟ](https://mni.wikipedia.org/wiki/\%EA\%AF\%89\%EA\%AF\%A5\%EA\%AF\%85\%EA\%AF\%A8)
\item [閩東語 / Mìng-dĕ̤ng-ngṳ̄](https://cdo.wikipedia.org/wiki/\%C3\%81k)
\item [Мокшень](https://mdf.wikipedia.org/wiki/\%D0\%AF\%D0\%BA\%D1\%81\%D1\%8F\%D1\%80\%D0\%B3\%D0\%B0)
\item [Монгол](https://mn.wikipedia.org/wiki/\%D0\%9D\%D1\%83\%D0\%B3\%D0\%B0\%D1\%81)
\item [မြန်မာဘာသာ](https://my.wikipedia.org/wiki/\%E1\%80\%98\%E1\%80\%B2)
\item [Nederlands](https://nl.wikipedia.org/wiki/Eenden)
\item [Nedersaksies](https://nds-nl.wikipedia.org/wiki/Ente)
\item [नेपाली](https://ne.wikipedia.org/wiki/\%E0\%A4\%B9\%E0\%A4\%BE\%E0\%A4\%81\%E0\%A4\%B8)
\item [नेपाल भाषा](https://new.wikipedia.org/wiki/\%E0\%A4\%B9\%E0\%A4\%81\%E0\%A4\%AF\%E0\%A5\%8D)
\item [日本語](https://ja.wikipedia.org/wiki/\%E3\%82\%AB\%E3\%83\%A2)
\item [Нохчийн](https://ce.wikipedia.org/wiki/\%D0\%91\%D0\%B5\%D0\%B4\%D0\%B0\%D1\%88)
\item [Norsk nynorsk](https://nn.wikipedia.org/wiki/And)
\item [Occitan](https://oc.wikipedia.org/wiki/Guit)
\item [Oromoo](https://om.wikipedia.org/wiki/Daakiyyee)
\item [ਪੰਜਾਬੀ](https://pa.wikipedia.org/wiki/\%E0\%A8\%AC\%E0\%A8\%A4\%E0\%A8\%96\%E0\%A8\%BC)
\item [Picard](https://pcd.wikipedia.org/wiki/Can\%C3\%A8rd)
\item [Plattdüütsch](https://nds.wikipedia.org/wiki/Aanten)
\item [Polski](https://pl.wikipedia.org/wiki/Kaczka)
\item [Português](https://pt.wikipedia.org/wiki/Pato)
\item [Qırımtatarca](https://crh.wikipedia.org/wiki/Papiy)
\item [Română](https://ro.wikipedia.org/wiki/Ra\%C8\%9B\%C4\%83)
\item [Русский](https://ru.wikipedia.org/wiki/\%D0\%A3\%D1\%82\%D0\%BA\%D0\%B8)
\item [Саха тыла](https://sah.wikipedia.org/wiki/\%D0\%9A\%D1\%83\%D1\%81\%D1\%82\%D0\%B0\%D1\%80)
\item [ᱥᱟᱱᱛᱟᱲᱤ](https://sat.wikipedia.org/wiki/\%E1\%B1\%9C\%E1\%B1\%AE\%E1\%B1\%B0\%E1\%B1\%AE)
\item [Sardu](https://sc.wikipedia.org/wiki/Anade)
\item [Scots](https://sco.wikipedia.org/wiki/Deuk)
\item [Seeltersk](https://stq.wikipedia.org/wiki/Oante)
\item [Shqip](https://sq.wikipedia.org/wiki/Rosa)
\item [Sicilianu](https://scn.wikipedia.org/wiki/P\%C3\%A0para\_(\%C3\%A0natra))
\item [සිංහල](https://si.wikipedia.org/wiki/\%E0\%B6\%AD\%E0\%B7\%8F\%E0\%B6\%BB\%E0\%B7\%8F\%E0\%B7\%80\%E0\%B7\%8F)
\item [Simple English](https://simple.wikipedia.org/wiki/Duck)
\item [سنڌي](https://sd.wikipedia.org/wiki/\%D8\%A8\%D8\%AF\%DA\%AA)
\item [کوردی](https://ckb.wikipedia.org/wiki/\%D9\%85\%D8\%B1\%D8\%A7\%D9\%88\%DB\%8C)
\item [Српски / srpski](https://sr.wikipedia.org/wiki/\%D0\%9F\%D0\%B0\%D1\%82\%D0\%BA\%D0\%B0)
\item [Srpskohrvatski / српскохрватски](https://sh.wikipedia.org/wiki/Patka)
\item [Sunda](https://su.wikipedia.org/wiki/Meri)
\item [Svenska](https://sv.wikipedia.org/wiki/Egentliga\_andf\%C3\%A5glar)
\item [Tagalog](https://tl.wikipedia.org/wiki/Bibi)
\item [தமிழ்](https://ta.wikipedia.org/wiki/\%E0\%AE\%B5\%E0\%AE\%BE\%E0\%AE\%A4\%E0\%AF\%8D\%E0\%AE\%A4\%E0\%AF\%81)
\item [Taqbaylit](https://kab.wikipedia.org/wiki/Ab\%E1\%B9\%9Bik)
\item [Татарча / tatarça](https://tt.wikipedia.org/wiki/\%D2\%AE\%D1\%80\%D0\%B4\%D3\%99\%D0\%BA\%D0\%BB\%D3\%99\%D1\%80)
\item [ไทย](https://th.wikipedia.org/wiki/\%E0\%B9\%80\%E0\%B8\%9B\%E0\%B9\%87\%E0\%B8\%94)
\item [Türkçe](https://tr.wikipedia.org/wiki/\%C3\%96rdek)
\item [Українська](https://uk.wikipedia.org/wiki/\%D0\%9A\%D0\%B0\%D1\%87\%D0\%BA\%D0\%B8)
\item [ئۇيغۇرچە / Uyghurche](https://ug.wikipedia.org/wiki/\%D8\%A6\%DB\%86\%D8\%B1\%D8\%AF\%DB\%95\%D9\%83)
\item [Vahcuengh](https://za.wikipedia.org/wiki/Bit\_(doenghduz))
\item [Tiếng Việt](https://vi.wikipedia.org/wiki/V\%E1\%BB\%8Bt)
\item [Walon](https://wa.wikipedia.org/wiki/Can\%C3\%A5rd)
\item [文言](https://zh-classical.wikipedia.org/wiki/\%E9\%B4\%A8)
\item [Winaray](https://war.wikipedia.org/wiki/Pato)
\item [吴语](https://wuu.wikipedia.org/wiki/\%E9\%B8\%AD)
\item [粵語](https://zh-yue.wikipedia.org/wiki/\%E9\%B4\%A8)
\item [Žemaitėška](https://bat-smg.wikipedia.org/wiki/P\%C4\%ABl\%C4\%97)
\item [中文](https://zh.wikipedia.org/wiki/\%E9\%B8\%AD)
\end{itemize}

[Edit links](https://www.wikidata.org/wiki/Special:EntityPage/Q3736439\#sitelinks-wikipedia)

\begin{itemize}
\item [Article](/wiki/Duck)
\item [Talk](/wiki/Talk:Duck)
\end{itemize}

- [ ] English

\begin{itemize}
\item [Read](/wiki/Duck)
\item [View source](/w/index.php?title=Duck\&action=edit)
\item [View history](/w/index.php?title=Duck\&action=history)
\end{itemize}

- [ ] Tools

Tools

move to sidebar

hide

Actions

\begin{itemize}
\item [Read](/wiki/Duck)
\item [View source](/w/index.php?title=Duck\&action=edit)
\item [View history](/w/index.php?title=Duck\&action=history)
\end{itemize}

General

\begin{itemize}
\item [What links here](/wiki/Special:WhatLinksHere/Duck)
\item [Related changes](/wiki/Special:RecentChangesLinked/Duck)
\item [Upload file](/wiki/Wikipedia:File\_Upload\_Wizard)
\item [Special pages](/wiki/Special:SpecialPages)
\item [Permanent link](/w/index.php?title=Duck\&oldid=1246843351)
\item [Page information](/w/index.php?title=Duck\&action=info)
\item [Cite this page](/w/index.php?title=Special:CiteThisPage\&page=Duck\&id=1246843351\&wpFormIdentifier=titleform)
\item [Get shortened URL](/w/index.php?title=Special:UrlShortener\&url=https\%3A\%2F\%2Fen.wikipedia.org\%2Fwiki\%2FDuck)
\item [Download QR code](/w/index.php?title=Special:QrCode\&url=https\%3A\%2F\%2Fen.wikipedia.org\%2Fwiki\%2FDuck)
\item [Wikidata item](https://www.wikidata.org/wiki/Special:EntityPage/Q3736439)
\end{itemize}

Print/export

\begin{itemize}
\item [Download as PDF](/w/index.php?title=Special:DownloadAsPdf\&page=Duck\&action=show-download-screen)
\item [Printable version](/w/index.php?title=Duck\&printable=yes)
\end{itemize}

In other projects

\begin{itemize}
\item [Wikimedia Commons](https://commons.wikimedia.org/wiki/Category:Ducks)
\item [Wikiquote](https://en.wikiquote.org/wiki/Duck)
\end{itemize}

Appearance

move to sidebar

hide

\begin{figure}[h]
% image
\caption{Page semi-protected}
\end{figure}

From Wikipedia, the free encyclopedia

(Redirected from [Duckling](/w/index.php?title=Duckling\&redirect=no) )

This article is about the bird. For duck as a food, see [Duck as food](/wiki/Duck\_as\_food) . For other uses, see [Duck (disambiguation)](/wiki/Duck\_(disambiguation)) .

"Duckling" redirects here. For other uses, see [Duckling (disambiguation)](/wiki/Duckling\_(disambiguation)) .

\begin{table}[h]
\begin{tabular}{|l|l|}
\hline
Duck & Duck \\ \hline
<!-- image --> & <!-- image --> \\ \hline
[Bufflehead](/wiki/Bufflehead) ( *Bucephala albeola* ) & [Bufflehead](/wiki/Bufflehead) ( *Bucephala albeola* ) \\ \hline
[Scientific classification](/wiki/Taxonomy\_(biology))  Edit this classification  <!-- image --> & [Scientific classification](/wiki/Taxonomy\_(biology))  Edit this classification  <!-- image --> \\ \hline
Domain: & [Eukaryota](/wiki/Eukaryote) \\ \hline
Kingdom: & [Animalia](/wiki/Animal) \\ \hline
Phylum: & [Chordata](/wiki/Chordate) \\ \hline
Class: & [Aves](/wiki/Bird) \\ \hline
Order: & [Anseriformes](/wiki/Anseriformes) \\ \hline
Superfamily: & [Anatoidea](/wiki/Anatoidea) \\ \hline
Family: & [Anatidae](/wiki/Anatidae) \\ \hline
Subfamilies & Subfamilies \\ \hline
See text & See text \\ \hline
\end{tabular}
\end{table}

**Duck** is the common name for numerous species of [waterfowl](/wiki/Waterfowl) in the [family](/wiki/Family\_(biology)) [Anatidae](/wiki/Anatidae) . Ducks are generally smaller and shorter-necked than [swans](/wiki/Swan) and [geese](/wiki/Goose) , which are members of the same family. Divided among several subfamilies, they are a [form taxon](/wiki/Form\_taxon) ; they do not represent a [monophyletic group](/wiki/Monophyletic\_group) (the group of all descendants of a single common ancestral species), since swans and geese are not considered ducks. Ducks are mostly [aquatic birds](/wiki/Aquatic\_bird) , and may be found in both fresh water and sea water.

Ducks are sometimes confused with several types of unrelated water birds with similar forms, such as [loons](/wiki/Loon) or divers, [grebes](/wiki/Grebe) , [gallinules](/wiki/Gallinule) and [coots](/wiki/Coot) .

\section{Etymology}

The word *duck* comes from [Old English](/wiki/Old\_English) *dūce* 'diver', a derivative of the verb * *dūcan* 'to duck, bend down low as if to get under something, or dive', because of the way many species in the [dabbling duck](/wiki/Dabbling\_duck) group feed by upending; compare with [Dutch](/wiki/Dutch\_language) *duiken* and [German](/wiki/German\_language) *tauchen* 'to dive'.

\begin{figure}[h]
% image
\caption{Pacific black duck displaying the characteristic upending "duck"}
\end{figure}

This word replaced Old English *ened* / *ænid* 'duck', possibly to avoid confusion with other words, such as *ende* 'end' with similar forms. Other Germanic languages still have similar words for *duck* , for example, Dutch *eend* , German *Ente* and [Norwegian](/wiki/Norwegian\_language) *and* . The word *ened* / *ænid* was inherited from [Proto-Indo-European](/wiki/Proto-Indo-European\_language) ; [cf.](/wiki/Cf.) [Latin](/wiki/Latin) *anas* "duck", [Lithuanian](/wiki/Lithuanian\_language) *ántis* 'duck', [Ancient Greek](/wiki/Ancient\_Greek\_language) νῆσσα / νῆττα ( *nēssa* / *nētta* ) 'duck', and [Sanskrit](/wiki/Sanskrit) *ātí* 'water bird', among others.

A duckling is a young duck in downy plumage [[ 1 ]](\#cite\_note-1) or baby duck, [[ 2 ]](\#cite\_note-2) but in the food trade a young domestic duck which has just reached adult size and bulk and its meat is still fully tender, is sometimes labelled as a duckling.

A male is called a [drake](https://en.wiktionary.org/wiki/drake) and the female is called a duck, or in [ornithology](/wiki/Ornithology) a hen. [[ 3 ]](\#cite\_note-3) [[ 4 ]](\#cite\_note-4)

\begin{figure}[h]
% image
\caption{Male mallard .}
\end{figure}

\begin{figure}[h]
% image
\caption{Wood ducks .}
\end{figure}

\section{Taxonomy}

All ducks belong to the [biological order](/wiki/Order\_(biology)) [Anseriformes](/wiki/Anseriformes) , a group that contains the ducks, geese and swans, as well as the [screamers](/wiki/Screamer) , and the [magpie goose](/wiki/Magpie\_goose) . [[ 5 ]](\#cite\_note-FOOTNOTECarboneras1992536-5) All except the screamers belong to the [biological family](/wiki/Family\_(biology)) [Anatidae](/wiki/Anatidae) . [[ 5 ]](\#cite\_note-FOOTNOTECarboneras1992536-5) Within the family, ducks are split into a variety of subfamilies and 'tribes'. The number and composition of these subfamilies and tribes is the cause of considerable disagreement among taxonomists. [[ 5 ]](\#cite\_note-FOOTNOTECarboneras1992536-5) Some base their decisions on [morphological characteristics](/wiki/Morphology\_(biology)) , others on shared behaviours or genetic studies. [[ 6 ]](\#cite\_note-FOOTNOTELivezey1986737–738-6) [[ 7 ]](\#cite\_note-FOOTNOTEMadsenMcHughde\_Kloet1988452-7) The number of suggested subfamilies containing ducks ranges from two to five. [[ 8 ]](\#cite\_note-FOOTNOTEDonne-GousséLaudetHänni2002353–354-8) [[ 9 ]](\#cite\_note-FOOTNOTECarboneras1992540-9) The significant level of [hybridisation](/wiki/Hybrid\_(biology)) that occurs among wild ducks complicates efforts to tease apart the relationships between various species. [[ 9 ]](\#cite\_note-FOOTNOTECarboneras1992540-9)

\begin{figure}[h]
% image
\caption{Mallard landing in approach}
\end{figure}

In most modern classifications, the so-called 'true ducks' belong to the subfamily Anatinae, which is further split into a varying number of tribes. [[ 10 ]](\#cite\_note-FOOTNOTEElphickDunningSibley2001191-10) The largest of these, the Anatini, contains the 'dabbling' or 'river' ducks - named for their method of feeding primarily at the surface of fresh water. [[ 11 ]](\#cite\_note-FOOTNOTEKear2005448-11) The 'diving ducks', also named for their primary feeding method, make up the tribe Aythyini. [[ 12 ]](\#cite\_note-FOOTNOTEKear2005622–623-12) The 'sea ducks' of the tribe Mergini are diving ducks which specialise on fish and shellfish and spend a majority of their lives in saltwater. [[ 13 ]](\#cite\_note-FOOTNOTEKear2005686-13) The tribe Oxyurini contains the 'stifftails', diving ducks notable for their small size and stiff, upright tails. [[ 14 ]](\#cite\_note-FOOTNOTEElphickDunningSibley2001193-14)

A number of other species called ducks are not considered to be 'true ducks', and are typically placed in other subfamilies or tribes. The [whistling ducks](/wiki/Whistling\_duck) are assigned either to a tribe (Dendrocygnini) in the subfamily Anatinae or the subfamily Anserinae, [[ 15 ]](\#cite\_note-FOOTNOTECarboneras1992537-15) or to their own subfamily (Dendrocygninae) or family (Dendrocyganidae). [[ 9 ]](\#cite\_note-FOOTNOTECarboneras1992540-9) [[ 16 ]](\#cite\_note-FOOTNOTEAmerican\_Ornithologists'\_Union1998xix-16) The [freckled duck](/wiki/Freckled\_duck) of Australia is either the sole member of the tribe Stictonettini in the subfamily Anserinae, [[ 15 ]](\#cite\_note-FOOTNOTECarboneras1992537-15) or in its own family, the Stictonettinae. [[ 9 ]](\#cite\_note-FOOTNOTECarboneras1992540-9) The [shelducks](/wiki/Shelduck) make up the tribe Tadornini in the family Anserinae in some classifications, [[ 15 ]](\#cite\_note-FOOTNOTECarboneras1992537-15) and their own subfamily, Tadorninae, in others, [[ 17 ]](\#cite\_note-FOOTNOTEAmerican\_Ornithologists'\_Union1998-17) while the [steamer ducks](/wiki/Steamer\_duck) are either placed in the family Anserinae in the tribe Tachyerini [[ 15 ]](\#cite\_note-FOOTNOTECarboneras1992537-15) or lumped with the shelducks in the tribe Tadorini. [[ 9 ]](\#cite\_note-FOOTNOTECarboneras1992540-9) The [perching ducks](/wiki/Perching\_duck) make up in the tribe Cairinini in the subfamily Anserinae in some classifications, while that tribe is eliminated in other classifications and its members assigned to the tribe Anatini. [[ 9 ]](\#cite\_note-FOOTNOTECarboneras1992540-9) The [torrent duck](/wiki/Torrent\_duck) is generally included in the subfamily Anserinae in the monotypic tribe Merganettini, [[ 15 ]](\#cite\_note-FOOTNOTECarboneras1992537-15) but is sometimes included in the tribe Tadornini. [[ 18 ]](\#cite\_note-FOOTNOTECarboneras1992538-18) The [pink-eared duck](/wiki/Pink-eared\_duck) is sometimes included as a true duck either in the tribe Anatini [[ 15 ]](\#cite\_note-FOOTNOTECarboneras1992537-15) or the tribe Malacorhynchini, [[ 19 ]](\#cite\_note-FOOTNOTEChristidisBoles200862-19) and other times is included with the shelducks in the tribe Tadornini. [[ 15 ]](\#cite\_note-FOOTNOTECarboneras1992537-15)

\section{Morphology}

\begin{figure}[h]
% image
\caption{Male Mandarin duck}
\end{figure}

The overall [body plan](/wiki/Body\_plan) of ducks is elongated and broad, and they are also relatively long-necked, albeit not as long-necked as the geese and swans. The body shape of diving ducks varies somewhat from this in being more rounded. The [bill](/wiki/Beak) is usually broad and contains serrated [pectens](/wiki/Pecten\_(biology)) , which are particularly well defined in the filter-feeding species. In the case of some fishing species the bill is long and strongly serrated. The scaled legs are strong and well developed, and generally set far back on the body, more so in the highly aquatic species. The wings are very strong and are generally short and pointed, and the [flight](/wiki/Bird\_flight) of ducks requires fast continuous strokes, requiring in turn strong wing muscles. Three species of [steamer duck](/wiki/Steamer\_duck) are almost flightless, however. Many species of duck are temporarily flightless while [moulting](/wiki/Moult) ; they seek out protected habitat with good food supplies during this period. This moult typically precedes [migration](/wiki/Bird\_migration) .

The drakes of northern species often have extravagant [plumage](/wiki/Plumage) , but that is [moulted](/wiki/Moult) in summer to give a more female-like appearance, the "eclipse" plumage. Southern resident species typically show less [sexual dimorphism](/wiki/Sexual\_dimorphism) , although there are exceptions such as the [paradise shelduck](/wiki/Paradise\_shelduck) of [New Zealand](/wiki/New\_Zealand) , which is both strikingly sexually dimorphic and in which the female's plumage is brighter than that of the male. The plumage of juvenile birds generally resembles that of the female. Female ducks have evolved to have a corkscrew shaped vagina to prevent rape.

\section{Distribution and habitat}

See also: [List of Anseriformes by population](/wiki/List\_of\_Anseriformes\_by\_population)

\begin{figure}[h]
% image
\caption{Flying steamer ducks in Ushuaia , Argentina}
\end{figure}

Ducks have a [cosmopolitan distribution](/wiki/Cosmopolitan\_distribution) , and are found on every continent except Antarctica. [[ 5 ]](\#cite\_note-FOOTNOTECarboneras1992536-5) Several species manage to live on subantarctic islands, including [South Georgia](/wiki/South\_Georgia\_and\_the\_South\_Sandwich\_Islands) and the [Auckland Islands](/wiki/Auckland\_Islands) . [[ 20 ]](\#cite\_note-FOOTNOTEShirihai2008239,\_245-20) Ducks have reached a number of isolated oceanic islands, including the [Hawaiian Islands](/wiki/Hawaiian\_Islands) , [Micronesia](/wiki/Micronesia) and the [Galápagos Islands](/wiki/Gal\%C3\%A1pagos\_Islands) , where they are often [vagrants](/wiki/Glossary\_of\_bird\_terms\#vagrants) and less often [residents](/wiki/Glossary\_of\_bird\_terms\#residents) . [[ 21 ]](\#cite\_note-FOOTNOTEPrattBrunerBerrett198798–107-21) [[ 22 ]](\#cite\_note-FOOTNOTEFitterFitterHosking200052–3-22) A handful are [endemic](/wiki/Endemic) to such far-flung islands. [[ 21 ]](\#cite\_note-FOOTNOTEPrattBrunerBerrett198798–107-21)

\begin{figure}[h]
% image
\caption{Female mallard in Cornwall , England}
\end{figure}

Some duck species, mainly those breeding in the temperate and Arctic Northern Hemisphere, are migratory; those in the tropics are generally not. Some ducks, particularly in Australia where rainfall is erratic, are nomadic, seeking out the temporary lakes and pools that form after localised heavy rain. [[ 23 ]](\#cite\_note-23)

\section{Behaviour}

\subsection{Feeding}

\begin{figure}[h]
% image
\caption{Pecten along the bill}
\end{figure}

Ducks eat food sources such as [grasses](/wiki/Poaceae) , aquatic plants, fish, insects, small amphibians, worms, and small [molluscs](/wiki/Mollusc) .

[Dabbling ducks](/wiki/Dabbling\_duck) feed on the surface of water or on land, or as deep as they can reach by up-ending without completely submerging. [[ 24 ]](\#cite\_note-24) Along the edge of the bill, there is a comb-like structure called a [pecten](/wiki/Pecten\_(biology)) . This strains the water squirting from the side of the bill and traps any food. The pecten is also used to preen feathers and to hold slippery food items.

[Diving ducks](/wiki/Diving\_duck) and [sea ducks](/wiki/Sea\_duck) forage deep underwater. To be able to submerge more easily, the diving ducks are heavier than dabbling ducks, and therefore have more difficulty taking off to fly.

A few specialized species such as the [mergansers](/wiki/Merganser) are adapted to catch and swallow large fish.

The others have the characteristic wide flat bill adapted to [dredging](/wiki/Dredging) -type jobs such as pulling up waterweed, pulling worms and small molluscs out of mud, searching for insect larvae, and bulk jobs such as dredging out, holding, turning head first, and swallowing a squirming frog. To avoid injury when digging into sediment it has no [cere](/wiki/Cere) , but the nostrils come out through hard horn.

[*The Guardian*](/wiki/The\_Guardian) published an article advising that ducks should not be fed with bread because it [damages the health of the ducks](/wiki/Angel\_wing) and pollutes waterways. [[ 25 ]](\#cite\_note-25)

\subsection{Breeding}

\begin{figure}[h]
% image
\caption{A Muscovy duckling}
\end{figure}

Ducks generally [only have one partner at a time](/wiki/Monogamy\_in\_animals) , although the partnership usually only lasts one year. [[ 26 ]](\#cite\_note-26) Larger species and the more sedentary species (like fast-river specialists) tend to have pair-bonds that last numerous years. [[ 27 ]](\#cite\_note-27) Most duck species breed once a year, choosing to do so in favourable conditions ( [spring](/wiki/Spring\_(season)) /summer or wet seasons). Ducks also tend to make a [nest](/wiki/Bird\_nest) before breeding, and, after hatching, lead their ducklings to water. Mother ducks are very caring and protective of their young, but may abandon some of their ducklings if they are physically stuck in an area they cannot get out of (such as nesting in an enclosed [courtyard](/wiki/Courtyard) ) or are not prospering due to genetic defects or sickness brought about by hypothermia, starvation, or disease. Ducklings can also be orphaned by inconsistent late hatching where a few eggs hatch after the mother has abandoned the nest and led her ducklings to water. [[ 28 ]](\#cite\_note-28)

\subsection{Communication}

Female [mallard](/wiki/Mallard) ducks (as well as several other species in the genus *Anas* , such as the [American](/wiki/American\_black\_duck) and [Pacific black ducks](/wiki/Pacific\_black\_duck) , [spot-billed duck](/wiki/Spot-billed\_duck) , [northern pintail](/wiki/Northern\_pintail) and [common teal](/wiki/Common\_teal) ) make the classic "quack" sound while males make a similar but raspier sound that is sometimes written as "breeeeze", [[ 29 ]](\#cite\_note-29) [ [*self-published source?*](/wiki/Wikipedia:Verifiability\#Self-published\_sources) ] but, despite widespread misconceptions, most species of duck do not "quack". [[ 30 ]](\#cite\_note-30) In general, ducks make a range of [calls](/wiki/Bird\_vocalisation) , including whistles, cooing, yodels and grunts. For example, the [scaup](/wiki/Scaup) - which are [diving ducks](/wiki/Diving\_duck) - make a noise like "scaup" (hence their name). Calls may be loud displaying calls or quieter contact calls.

A common [urban legend](/wiki/Urban\_legend) claims that duck quacks do not echo; however, this has been proven to be false. This myth was first debunked by the Acoustics Research Centre at the [University of Salford](/wiki/University\_of\_Salford) in 2003 as part of the [British Association](/wiki/British\_Association) 's Festival of Science. [[ 31 ]](\#cite\_note-31) It was also debunked in [one of the earlier episodes](/wiki/MythBusters\_(2003\_season)\#Does\_a\_Duck's\_Quack\_Echo?) of the popular Discovery Channel television show [*MythBusters*](/wiki/MythBusters) . [[ 32 ]](\#cite\_note-32)

\subsection{Predators}

\begin{figure}[h]
% image
\caption{Ringed teal}
\end{figure}

Ducks have many predators. Ducklings are particularly vulnerable, since their inability to fly makes them easy prey not only for predatory birds but also for large fish like [pike](/wiki/Esox) , [crocodilians](/wiki/Crocodilia) , predatory [testudines](/wiki/Testudines) such as the [alligator snapping turtle](/wiki/Alligator\_snapping\_turtle) , and other aquatic hunters, including fish-eating birds such as [herons](/wiki/Heron) . Ducks' nests are raided by land-based predators, and brooding females may be caught unaware on the nest by mammals, such as [foxes](/wiki/Fox) , or large birds, such as [hawks](/wiki/Hawk) or [owls](/wiki/Owl) .

Adult ducks are fast fliers, but may be caught on the water by large aquatic predators including big fish such as the North American [muskie](/wiki/Muskellunge) and the European [pike](/wiki/Esox) . In flight, ducks are safe from all but a few predators such as humans and the [peregrine falcon](/wiki/Peregrine\_falcon) , which uses its speed and strength to catch ducks.

\section{Relationship with humans}

\subsection{Hunting}

Main article: [Waterfowl hunting](/wiki/Waterfowl\_hunting)

Humans have hunted ducks since prehistoric times. Excavations of [middens](/wiki/Midden) in California dating to 7800 - 6400 [BP](/wiki/Before\_present) have turned up bones of ducks, including at least one now-extinct flightless species. [[ 33 ]](\#cite\_note-FOOTNOTEErlandson1994171-33) Ducks were captured in "significant numbers" by [Holocene](/wiki/Holocene) inhabitants of the lower [Ohio River](/wiki/Ohio\_River) valley, suggesting they took advantage of the seasonal bounty provided by migrating waterfowl. [[ 34 ]](\#cite\_note-FOOTNOTEJeffries2008168,\_243-34) Neolithic hunters in locations as far apart as the Caribbean, [[ 35 ]](\#cite\_note-FOOTNOTESued-Badillo200365-35) Scandinavia, [[ 36 ]](\#cite\_note-FOOTNOTEThorpe199668-36) Egypt, [[ 37 ]](\#cite\_note-FOOTNOTEMaisels199942-37) Switzerland, [[ 38 ]](\#cite\_note-FOOTNOTERau1876133-38) and China relied on ducks as a source of protein for some or all of the year. [[ 39 ]](\#cite\_note-FOOTNOTEHigman201223-39) Archeological evidence shows that [Māori people](/wiki/M\%C4\%81ori\_people) in New Zealand hunted the flightless [Finsch's duck](/wiki/Finsch\%27s\_duck) , possibly to extinction, though rat predation may also have contributed to its fate. [[ 40 ]](\#cite\_note-FOOTNOTEHume201253-40) A similar end awaited the [Chatham duck](/wiki/Chatham\_duck) , a species with reduced flying capabilities which went extinct shortly after its island was colonised by Polynesian settlers. [[ 41 ]](\#cite\_note-FOOTNOTEHume201252-41) It is probable that duck eggs were gathered by Neolithic hunter-gathers as well, though hard evidence of this is uncommon. [[ 35 ]](\#cite\_note-FOOTNOTESued-Badillo200365-35) [[ 42 ]](\#cite\_note-FOOTNOTEFieldhouse2002167-42)

In many areas, wild ducks (including ducks farmed and released into the wild) are hunted for food or sport, [[ 43 ]](\#cite\_note-43) by shooting, or by being trapped using [duck decoys](/wiki/Duck\_decoy\_(structure)) . Because an idle floating duck or a duck squatting on land cannot react to fly or move quickly, "a sitting duck" has come to mean "an easy target". These ducks may be [contaminated by pollutants](/wiki/Duck\_(food)\#Pollution) such as [PCBs](/wiki/Polychlorinated\_biphenyl) . [[ 44 ]](\#cite\_note-44)

\subsection{Domestication}

Main article: [Domestic duck](/wiki/Domestic\_duck)

\begin{figure}[h]
% image
\caption{Indian Runner ducks , a common breed of domestic ducks}
\end{figure}

Ducks have many economic uses, being farmed for their meat, eggs, and feathers (particularly their [down](/wiki/Down\_feather) ). Approximately 3 billion ducks are slaughtered each year for meat worldwide. [[ 45 ]](\#cite\_note-45) They are also kept and bred by aviculturists and often displayed in zoos. Almost all the varieties of domestic ducks are descended from the [mallard](/wiki/Mallard) ( *Anas platyrhynchos* ), apart from the [Muscovy duck](/wiki/Muscovy\_duck) ( *Cairina moschata* ). [[ 46 ]](\#cite\_note-46) [[ 47 ]](\#cite\_note-47) The [Call duck](/wiki/Call\_duck) is another example of a domestic duck breed. Its name comes from its original use established by hunters, as a decoy to attract wild mallards from the sky, into traps set for them on the ground. The call duck is the world's smallest domestic duck breed, as it weighs less than 1 kg (2.2 lb). [[ 48 ]](\#cite\_note-48)

\subsection{Heraldry}

\begin{figure}[h]
% image
\caption{Three black-colored ducks in the coat of arms of Maaninka [ 49 ]}
\end{figure}

Ducks appear on several [coats of arms](/wiki/Coats\_of\_arms) , including the coat of arms of [Lubāna](/wiki/Lub\%C4\%81na) ( [Latvia](/wiki/Latvia) ) [[ 50 ]](\#cite\_note-50) and the coat of arms of [Föglö](/wiki/F\%C3\%B6gl\%C3\%B6) ( [Åland](/wiki/\%C3\%85land) ). [[ 51 ]](\#cite\_note-51)

\subsection{Cultural references}

In 2002, psychologist [Richard Wiseman](/wiki/Richard\_Wiseman) and colleagues at the [University of Hertfordshire](/wiki/University\_of\_Hertfordshire) , [UK](/wiki/UK) , finished a year-long [LaughLab](/wiki/LaughLab) experiment, concluding that of all animals, ducks attract the most humor and silliness; he said, "If you're going to tell a joke involving an animal, make it a duck." [[ 52 ]](\#cite\_note-52) The word "duck" may have become an [inherently funny word](/wiki/Inherently\_funny\_word) in many languages, possibly because ducks are seen as silly in their looks or behavior. Of the many [ducks in fiction](/wiki/List\_of\_fictional\_ducks) , many are cartoon characters, such as [Walt Disney](/wiki/The\_Walt\_Disney\_Company) 's [Donald Duck](/wiki/Donald\_Duck) , and [Warner Bros.](/wiki/Warner\_Bros.) ' [Daffy Duck](/wiki/Daffy\_Duck) . [Howard the Duck](/wiki/Howard\_the\_Duck) started as a comic book character in 1973 [[ 53 ]](\#cite\_note-53) [[ 54 ]](\#cite\_note-54) and was made into a [movie](/wiki/Howard\_the\_Duck\_(film)) in 1986.

The 1992 Disney film [*The Mighty Ducks*](/wiki/The\_Mighty\_Ducks\_(film)) , starring [Emilio Estevez](/wiki/Emilio\_Estevez) , chose the duck as the mascot for the fictional youth hockey team who are protagonists of the movie, based on the duck being described as a fierce fighter. This led to the duck becoming the nickname and mascot for the eventual [National Hockey League](/wiki/National\_Hockey\_League) professional team of the [Anaheim Ducks](/wiki/Anaheim\_Ducks) , who were founded with the name the Mighty Ducks of Anaheim. [ [*citation needed*](/wiki/Wikipedia:Citation\_needed) ] The duck is also the nickname of the [University of Oregon](/wiki/University\_of\_Oregon) sports teams as well as the [Long Island Ducks](/wiki/Long\_Island\_Ducks) minor league [baseball](/wiki/Baseball) team. [[ 55 ]](\#cite\_note-55)

\section{See also}

\begin{itemize}
\item [Birds portal](/wiki/Portal:Birds)
<!-- image -->
\end{itemize}

\begin{itemize}
\item [Domestic duck](/wiki/Domestic\_duck)
\item [Duck as food](/wiki/Duck\_as\_food)
\item [Duck test](/wiki/Duck\_test)
\item [Duck breeds](/wiki/List\_of\_duck\_breeds)
\item [Fictional ducks](/wiki/List\_of\_fictional\_ducks)
\item [Rubber duck](/wiki/Rubber\_duck)
\end{itemize}

\section{Notes}

\subsection{Citations}

\begin{enumerate}
\item [**\textasciicircum{}**](\#cite\_ref-1) ["Duckling"](http://dictionary.reference.com/browse/duckling) . *The American Heritage Dictionary of the English Language, Fourth Edition* . Houghton Mifflin Company. 2006 . Retrieved 2015-05-22 .
\item [**\textasciicircum{}**](\#cite\_ref-2) ["Duckling"](http://dictionary.reference.com/browse/duckling) . *Kernerman English Multilingual Dictionary (Beta Version)* . K. Dictionaries Ltd. 2000-2006 . Retrieved 2015-05-22 .
\item [**\textasciicircum{}**](\#cite\_ref-3) Dohner, Janet Vorwald (2001). [*The Encyclopedia of Historic and Endangered Livestock and Poultry Breeds*](https://books.google.com/books?id=WJCTL\_mC5w4C\&q=male+duck+is+called+a+drake+and+the+female+is+called+a+duck\&pg=PA457) . Yale University Press. [ISBN](/wiki/ISBN\_(identifier)) [978-0300138139](/wiki/Special:BookSources/978-0300138139) .
\item [**\textasciicircum{}**](\#cite\_ref-4) Visca, Curt; Visca, Kelley (2003). [*How to Draw Cartoon Birds*](https://books.google.com/books?id=VqSquCLNrZcC\&q=male+duck+is+called+a+drake+and+the+female+is+called+a+duck+\%28or+hen\%29\&pg=PA16) . The Rosen Publishing Group. [ISBN](/wiki/ISBN\_(identifier)) [9780823961566](/wiki/Special:BookSources/9780823961566) .
\item \textasciicircum{} [***a***](\#cite\_ref-FOOTNOTECarboneras1992536\_5-0) [***b***](\#cite\_ref-FOOTNOTECarboneras1992536\_5-1) [***c***](\#cite\_ref-FOOTNOTECarboneras1992536\_5-2) [***d***](\#cite\_ref-FOOTNOTECarboneras1992536\_5-3) [Carboneras 1992](\#CITEREFCarboneras1992) , p. 536.
\item [**\textasciicircum{}**](\#cite\_ref-FOOTNOTELivezey1986737–738\_6-0) [Livezey 1986](\#CITEREFLivezey1986) , pp. 737-738.
\item [**\textasciicircum{}**](\#cite\_ref-FOOTNOTEMadsenMcHughde\_Kloet1988452\_7-0) [Madsen, McHugh \& de Kloet 1988](\#CITEREFMadsenMcHughde\_Kloet1988) , p. 452.
\item [**\textasciicircum{}**](\#cite\_ref-FOOTNOTEDonne-GousséLaudetHänni2002353–354\_8-0) [Donne-Goussé, Laudet \& Hänni 2002](\#CITEREFDonne-GousséLaudetHänni2002) , pp. 353-354.
\item \textasciicircum{} [***a***](\#cite\_ref-FOOTNOTECarboneras1992540\_9-0) [***b***](\#cite\_ref-FOOTNOTECarboneras1992540\_9-1) [***c***](\#cite\_ref-FOOTNOTECarboneras1992540\_9-2) [***d***](\#cite\_ref-FOOTNOTECarboneras1992540\_9-3) [***e***](\#cite\_ref-FOOTNOTECarboneras1992540\_9-4) [***f***](\#cite\_ref-FOOTNOTECarboneras1992540\_9-5) [Carboneras 1992](\#CITEREFCarboneras1992) , p. 540.
\item [**\textasciicircum{}**](\#cite\_ref-FOOTNOTEElphickDunningSibley2001191\_10-0) [Elphick, Dunning \& Sibley 2001](\#CITEREFElphickDunningSibley2001) , p. 191.
\item [**\textasciicircum{}**](\#cite\_ref-FOOTNOTEKear2005448\_11-0) [Kear 2005](\#CITEREFKear2005) , p. 448.
\item [**\textasciicircum{}**](\#cite\_ref-FOOTNOTEKear2005622–623\_12-0) [Kear 2005](\#CITEREFKear2005) , p. 622-623.
\item [**\textasciicircum{}**](\#cite\_ref-FOOTNOTEKear2005686\_13-0) [Kear 2005](\#CITEREFKear2005) , p. 686.
\item [**\textasciicircum{}**](\#cite\_ref-FOOTNOTEElphickDunningSibley2001193\_14-0) [Elphick, Dunning \& Sibley 2001](\#CITEREFElphickDunningSibley2001) , p. 193.
\item \textasciicircum{} [***a***](\#cite\_ref-FOOTNOTECarboneras1992537\_15-0) [***b***](\#cite\_ref-FOOTNOTECarboneras1992537\_15-1) [***c***](\#cite\_ref-FOOTNOTECarboneras1992537\_15-2) [***d***](\#cite\_ref-FOOTNOTECarboneras1992537\_15-3) [***e***](\#cite\_ref-FOOTNOTECarboneras1992537\_15-4) [***f***](\#cite\_ref-FOOTNOTECarboneras1992537\_15-5) [***g***](\#cite\_ref-FOOTNOTECarboneras1992537\_15-6) [Carboneras 1992](\#CITEREFCarboneras1992) , p. 537.
\item [**\textasciicircum{}**](\#cite\_ref-FOOTNOTEAmerican\_Ornithologists'\_Union1998xix\_16-0) [American Ornithologists' Union 1998](\#CITEREFAmerican\_Ornithologists'\_Union1998) , p. xix.
\item [**\textasciicircum{}**](\#cite\_ref-FOOTNOTEAmerican\_Ornithologists'\_Union1998\_17-0) [American Ornithologists' Union 1998](\#CITEREFAmerican\_Ornithologists'\_Union1998) .
\item [**\textasciicircum{}**](\#cite\_ref-FOOTNOTECarboneras1992538\_18-0) [Carboneras 1992](\#CITEREFCarboneras1992) , p. 538.
\item [**\textasciicircum{}**](\#cite\_ref-FOOTNOTEChristidisBoles200862\_19-0) [Christidis \& Boles 2008](\#CITEREFChristidisBoles2008) , p. 62.
\item [**\textasciicircum{}**](\#cite\_ref-FOOTNOTEShirihai2008239,\_245\_20-0) [Shirihai 2008](\#CITEREFShirihai2008) , pp. 239, 245.
\item \textasciicircum{} [***a***](\#cite\_ref-FOOTNOTEPrattBrunerBerrett198798–107\_21-0) [***b***](\#cite\_ref-FOOTNOTEPrattBrunerBerrett198798–107\_21-1) [Pratt, Bruner \& Berrett 1987](\#CITEREFPrattBrunerBerrett1987) , pp. 98-107.
\item [**\textasciicircum{}**](\#cite\_ref-FOOTNOTEFitterFitterHosking200052–3\_22-0) [Fitter, Fitter \& Hosking 2000](\#CITEREFFitterFitterHosking2000) , pp. 52-3.
\item [**\textasciicircum{}**](\#cite\_ref-23) ["Pacific Black Duck"](http://www.wiresnr.org/pacificblackduck.html) . *www.wiresnr.org* . Retrieved 2018-04-27 .
\item [**\textasciicircum{}**](\#cite\_ref-24) Ogden, Evans. ["Dabbling Ducks"](https://www.sfu.ca/biology/wildberg/species/dabbducks.html) . CWE . Retrieved 2006-11-02 .
\item [**\textasciicircum{}**](\#cite\_ref-25) Karl Mathiesen (16 March 2015). ["Don't feed the ducks bread, say conservationists"](https://www.theguardian.com/environment/2015/mar/16/dont-feed-the-ducks-bread-say-conservationists) . *The Guardian* . Retrieved 13 November 2016 .
\item [**\textasciicircum{}**](\#cite\_ref-26) Rohwer, Frank C.; Anderson, Michael G. (1988). "Female-Biased Philopatry, Monogamy, and the Timing of Pair Formation in Migratory Waterfowl". *Current Ornithology* . pp. 187-221. [doi](/wiki/Doi\_(identifier)) : [10.1007/978-1-4615-6787-5\_4](https://doi.org/10.1007\%2F978-1-4615-6787-5\_4) . [ISBN](/wiki/ISBN\_(identifier)) [978-1-4615-6789-9](/wiki/Special:BookSources/978-1-4615-6789-9) .
\item [**\textasciicircum{}**](\#cite\_ref-27) Smith, Cyndi M.; Cooke, Fred; Robertson, Gregory J.; Goudie, R. Ian; Boyd, W. Sean (2000). ["Long-Term Pair Bonds in Harlequin Ducks"](https://doi.org/10.1093\%2Fcondor\%2F102.1.201) . *The Condor* . **102** (1): 201-205. [doi](/wiki/Doi\_(identifier)) : [10.1093/condor/102.1.201](https://doi.org/10.1093\%2Fcondor\%2F102.1.201) . [hdl](/wiki/Hdl\_(identifier)) : [10315/13797](https://hdl.handle.net/10315\%2F13797) .
\item [**\textasciicircum{}**](\#cite\_ref-28) ["If You Find An Orphaned Duckling - Wildlife Rehabber"](https://web.archive.org/web/20180923152911/http://wildliferehabber.com/content/if-you-find-duckling) . *wildliferehabber.com* . Archived from [the original](https://wildliferehabber.com/content/if-you-find-duckling) on 2018-09-23 . Retrieved 2018-12-22 .
\item [**\textasciicircum{}**](\#cite\_ref-29) Carver, Heather (2011). [*The Duck Bible*](https://books.google.com/books?id=VGofAwAAQBAJ\&q=mallard+sound+deep+and+raspy\&pg=PA39) . Lulu.com. [ISBN](/wiki/ISBN\_(identifier)) [9780557901562](/wiki/Special:BookSources/9780557901562) . [ [*self-published source*](/wiki/Wikipedia:Verifiability\#Self-published\_sources) ]
\item [**\textasciicircum{}**](\#cite\_ref-30) Titlow, Budd (2013-09-03). [*Bird Brains: Inside the Strange Minds of Our Fine Feathered Friends*](https://books.google.com/books?id=fXJBBAAAQBAJ\&q=Females+of+most+dabbling+ducks+make+the+classic+\%22quack\%22+sound+but+most+ducks+don\%27t+quack\&pg=PA123) . Rowman \& Littlefield. [ISBN](/wiki/ISBN\_(identifier)) [9780762797707](/wiki/Special:BookSources/9780762797707) .
\item [**\textasciicircum{}**](\#cite\_ref-31) Amos, Jonathan (2003-09-08). ["Sound science is quackers"](http://news.bbc.co.uk/2/hi/science/nature/3086890.stm) . *BBC News* . Retrieved 2006-11-02 .
\item [**\textasciicircum{}**](\#cite\_ref-32) ["Mythbusters Episode 8"](http://mythbustersresults.com/episode8) . 12 December 2003.
\item [**\textasciicircum{}**](\#cite\_ref-FOOTNOTEErlandson1994171\_33-0) [Erlandson 1994](\#CITEREFErlandson1994) , p. 171.
\item [**\textasciicircum{}**](\#cite\_ref-FOOTNOTEJeffries2008168,\_243\_34-0) [Jeffries 2008](\#CITEREFJeffries2008) , pp. 168, 243.
\item \textasciicircum{} [***a***](\#cite\_ref-FOOTNOTESued-Badillo200365\_35-0) [***b***](\#cite\_ref-FOOTNOTESued-Badillo200365\_35-1) [Sued-Badillo 2003](\#CITEREFSued-Badillo2003) , p. 65.
\item [**\textasciicircum{}**](\#cite\_ref-FOOTNOTEThorpe199668\_36-0) [Thorpe 1996](\#CITEREFThorpe1996) , p. 68.
\item [**\textasciicircum{}**](\#cite\_ref-FOOTNOTEMaisels199942\_37-0) [Maisels 1999](\#CITEREFMaisels1999) , p. 42.
\item [**\textasciicircum{}**](\#cite\_ref-FOOTNOTERau1876133\_38-0) [Rau 1876](\#CITEREFRau1876) , p. 133.
\item [**\textasciicircum{}**](\#cite\_ref-FOOTNOTEHigman201223\_39-0) [Higman 2012](\#CITEREFHigman2012) , p. 23.
\item [**\textasciicircum{}**](\#cite\_ref-FOOTNOTEHume201253\_40-0) [Hume 2012](\#CITEREFHume2012) , p. 53.
\item [**\textasciicircum{}**](\#cite\_ref-FOOTNOTEHume201252\_41-0) [Hume 2012](\#CITEREFHume2012) , p. 52.
\item [**\textasciicircum{}**](\#cite\_ref-FOOTNOTEFieldhouse2002167\_42-0) [Fieldhouse 2002](\#CITEREFFieldhouse2002) , p. 167.
\item [**\textasciicircum{}**](\#cite\_ref-43) Livingston, A. D. (1998-01-01). [*Guide to Edible Plants and Animals*](https://books.google.com/books?id=NViSMffyaSgC\&q=\%C2\%A0\%C2\%A0In+many+areas,+wild+ducks+of+various+species+are+hunted+for+food+or+sport) . Wordsworth Editions, Limited. [ISBN](/wiki/ISBN\_(identifier)) [9781853263774](/wiki/Special:BookSources/9781853263774) .
\item [**\textasciicircum{}**](\#cite\_ref-44) ["Study plan for waterfowl injury assessment: Determining PCB concentrations in Hudson river resident waterfowl"](https://www.dec.ny.gov/docs/fish\_marine\_pdf/wfp09a.pdf) (PDF) . *New York State Department of Environmental Conservation* . US Department of Commerce. December 2008. p. 3. [Archived](https://ghostarchive.org/archive/20221009/https://www.dec.ny.gov/docs/fish\_marine\_pdf/wfp09a.pdf) (PDF) from the original on 2022-10-09 . Retrieved 2 July 2019 .
\item [**\textasciicircum{}**](\#cite\_ref-45) ["FAOSTAT"](http://www.fao.org/faostat/en/\#data/QL) . *www.fao.org* . Retrieved 2019-10-25 .
\item [**\textasciicircum{}**](\#cite\_ref-46) ["Anas platyrhynchos, Domestic Duck; DigiMorph Staff - The University of Texas at Austin"](http://digimorph.org/specimens/anas\_platyrhynchos/skull/) . Digimorph.org . Retrieved 2012-12-23 .
\item [**\textasciicircum{}**](\#cite\_ref-47) Sy Montgomery. ["Mallard; Encyclopædia Britannica"](https://www.britannica.com/eb/topic-360302/mallard) . *Britannica.com* . Retrieved 2012-12-23 .
\item [**\textasciicircum{}**](\#cite\_ref-48) Glenday, Craig (2014). [*Guinness World Records*](https://archive.org/details/guinnessworldrec0000unse\_r3e7/page/135) . Guinness World Records Limited. pp. [135](https://archive.org/details/guinnessworldrec0000unse\_r3e7/page/135) . [ISBN](/wiki/ISBN\_(identifier)) [978-1-908843-15-9](/wiki/Special:BookSources/978-1-908843-15-9) .
\item [**\textasciicircum{}**](\#cite\_ref-49) *Suomen kunnallisvaakunat* (in Finnish). Suomen Kunnallisliitto. 1982. p. 147. [ISBN](/wiki/ISBN\_(identifier)) [951-773-085-3](/wiki/Special:BookSources/951-773-085-3) .
\item [**\textasciicircum{}**](\#cite\_ref-50) ["Lubānas simbolika"](http://www.lubana.lv/index.php/lv/homepage/lubanas-pilseta-2) (in Latvian) . Retrieved September 9, 2021 .
\item [**\textasciicircum{}**](\#cite\_ref-51) ["Föglö"](http://digi.narc.fi/digi/view.ka?kuid=1738595) (in Swedish) . Retrieved September 9, 2021 .
\item [**\textasciicircum{}**](\#cite\_ref-52) Young, Emma. ["World's funniest joke revealed"](https://www.newscientist.com/article/dn2876-worlds-funniest-joke-revealed/) . *New Scientist* . Retrieved 7 January 2019 .
\item [**\textasciicircum{}**](\#cite\_ref-53) ["Howard the Duck (character)"](http://www.comics.org/character/name/Howard\%20the\%20Duck/sort/chrono/) . [*Grand Comics Database*](/wiki/Grand\_Comics\_Database) .
\item [**\textasciicircum{}**](\#cite\_ref-54) [Sanderson, Peter](/wiki/Peter\_Sanderson) ; Gilbert, Laura (2008). "1970s". *Marvel Chronicle A Year by Year History* . London, United Kingdom: [Dorling Kindersley](/wiki/Dorling\_Kindersley) . p. 161. [ISBN](/wiki/ISBN\_(identifier)) [978-0756641238](/wiki/Special:BookSources/978-0756641238) . December saw the debut of the cigar-smoking Howard the Duck. In this story by writer Steve Gerber and artist Val Mayerik, various beings from different realities had begun turning up in the Man-Thing's Florida swamp, including this bad-tempered talking duck.
\item [**\textasciicircum{}**](\#cite\_ref-55) ["The Duck"](https://goducks.com/sports/2003/8/28/153778.aspx) . *University of Oregon Athletics* . Retrieved 2022-01-20 .
\end{enumerate}

\subsection{Sources}

\begin{itemize}
\item American Ornithologists' Union (1998). [*Checklist of North American Birds*](https://americanornithology.org/wp-content/uploads/2019/07/AOSChecklistTin-Falcon.pdf) (PDF) . Washington, DC: American Ornithologists' Union. [ISBN](/wiki/ISBN\_(identifier)) [978-1-891276-00-2](/wiki/Special:BookSources/978-1-891276-00-2) . [Archived](https://ghostarchive.org/archive/20221009/https://americanornithology.org/wp-content/uploads/2019/07/AOSChecklistTin-Falcon.pdf) (PDF) from the original on 2022-10-09.
\item Carboneras, Carlos (1992). del Hoyo, Josep; Elliott, Andrew; Sargatal, Jordi (eds.). *Handbook of the Birds of the World* . Vol. 1: Ostrich to Ducks. Barcelona: Lynx Edicions. [ISBN](/wiki/ISBN\_(identifier)) [978-84-87334-10-8](/wiki/Special:BookSources/978-84-87334-10-8) .
\item Christidis, Les; Boles, Walter E., eds. (2008). *Systematics and Taxonomy of Australian Birds* . Collingwood, VIC: Csiro Publishing. [ISBN](/wiki/ISBN\_(identifier)) [978-0-643-06511-6](/wiki/Special:BookSources/978-0-643-06511-6) .
\item Donne-Goussé, Carole; Laudet, Vincent; Hänni, Catherine (July 2002). "A molecular phylogeny of Anseriformes based on mitochondrial DNA analysis". *Molecular Phylogenetics and Evolution* . **23** (3): 339-356. [Bibcode](/wiki/Bibcode\_(identifier)) : [2002MolPE..23..339D](https://ui.adsabs.harvard.edu/abs/2002MolPE..23..339D) . [doi](/wiki/Doi\_(identifier)) : [10.1016/S1055-7903(02)00019-2](https://doi.org/10.1016\%2FS1055-7903\%2802\%2900019-2) . [PMID](/wiki/PMID\_(identifier)) [12099792](https://pubmed.ncbi.nlm.nih.gov/12099792) .
\item Elphick, Chris; Dunning, John B. Jr.; Sibley, David, eds. (2001). *The Sibley Guide to Bird Life and Behaviour* . London: Christopher Helm. [ISBN](/wiki/ISBN\_(identifier)) [978-0-7136-6250-4](/wiki/Special:BookSources/978-0-7136-6250-4) .
\item Erlandson, Jon M. (1994). [*Early Hunter-Gatherers of the California Coast*](https://books.google.com/books?id=nGTaBwAAQBAJ\&pg=171) . New York, NY: Springer Science \& Business Media. [ISBN](/wiki/ISBN\_(identifier)) [978-1-4419-3231-0](/wiki/Special:BookSources/978-1-4419-3231-0) .
\item Fieldhouse, Paul (2002). [*Food, Feasts, and Faith: An Encyclopedia of Food Culture in World Religions*](https://books.google.com/books?id=P-FqDgAAQBAJ\&pg=PA167) . Vol. I: A-K. Santa Barbara: ABC-CLIO. [ISBN](/wiki/ISBN\_(identifier)) [978-1-61069-412-4](/wiki/Special:BookSources/978-1-61069-412-4) .
\item Fitter, Julian; Fitter, Daniel; Hosking, David (2000). *Wildlife of the Galápagos* . Princeton, NJ: Princeton University Press. [ISBN](/wiki/ISBN\_(identifier)) [978-0-691-10295-5](/wiki/Special:BookSources/978-0-691-10295-5) .
\item Higman, B. W. (2012). [*How Food Made History*](https://books.google.com/books?id=YIUoz98yMvgC\&pg=RA1-PA1801) . Chichester, UK: John Wiley \& Sons. [ISBN](/wiki/ISBN\_(identifier)) [978-1-4051-8947-7](/wiki/Special:BookSources/978-1-4051-8947-7) .
\item Hume, Julian H. (2012). [*Extinct Birds*](https://books.google.com/books?id=40sxDwAAQBAJ\&pg=PA53) . London: Christopher Helm. [ISBN](/wiki/ISBN\_(identifier)) [978-1-4729-3744-5](/wiki/Special:BookSources/978-1-4729-3744-5) .
\item Jeffries, Richard (2008). [*Holocene Hunter-Gatherers of the Lower Ohio River Valley*](https://archive.org/details/holocenehunterga0000jeff/mode/2up) . Tuscaloosa: University of Alabama Press. [ISBN](/wiki/ISBN\_(identifier)) [978-0-8173-1658-7](/wiki/Special:BookSources/978-0-8173-1658-7) .
\item Kear, Janet, ed. (2005). *Ducks, Geese and Swans: Species Accounts (* Cairina *to* Mergus *)* . Bird Families of the World. Oxford: Oxford University Press. [ISBN](/wiki/ISBN\_(identifier)) [978-0-19-861009-0](/wiki/Special:BookSources/978-0-19-861009-0) .
\item Livezey, Bradley C. (October 1986). ["A phylogenetic analysis of recent Anseriform genera using morphological characters"](https://sora.unm.edu/sites/default/files/journals/auk/v103n04/p0737-p0754.pdf) (PDF) . *The Auk* . **103** (4): 737-754. [doi](/wiki/Doi\_(identifier)) : [10.1093/auk/103.4.737](https://doi.org/10.1093\%2Fauk\%2F103.4.737) . [Archived](https://ghostarchive.org/archive/20221009/https://sora.unm.edu/sites/default/files/journals/auk/v103n04/p0737-p0754.pdf) (PDF) from the original on 2022-10-09.
\item Madsen, Cort S.; McHugh, Kevin P.; de Kloet, Siwo R. (July 1988). ["A partial classification of waterfowl (Anatidae) based on single-copy DNA"](https://sora.unm.edu/sites/default/files/journals/auk/v105n03/p0452-p0459.pdf) (PDF) . *The Auk* . **105** (3): 452-459. [doi](/wiki/Doi\_(identifier)) : [10.1093/auk/105.3.452](https://doi.org/10.1093\%2Fauk\%2F105.3.452) . [Archived](https://ghostarchive.org/archive/20221009/https://sora.unm.edu/sites/default/files/journals/auk/v105n03/p0452-p0459.pdf) (PDF) from the original on 2022-10-09.
\item Maisels, Charles Keith (1999). [*Early Civilizations of the Old World*](https://books.google.com/books?id=I2dgI2ijww8C\&pg=PA42) . London: Routledge. [ISBN](/wiki/ISBN\_(identifier)) [978-0-415-10975-8](/wiki/Special:BookSources/978-0-415-10975-8) .
\item Pratt, H. Douglas; Bruner, Phillip L.; Berrett, Delwyn G. (1987). *A Field Guide to the Birds of Hawaii and the Tropical Pacific* . Princeton, NJ: Princeton University Press. [ISBN](/wiki/ISBN\_(identifier)) [0-691-02399-9](/wiki/Special:BookSources/0-691-02399-9) .
\item Rau, Charles (1876). [*Early Man in Europe*](https://books.google.com/books?id=9XBgAAAAIAAJ\&pg=133) . New York: Harper \& Brothers. [LCCN](/wiki/LCCN\_(identifier)) [05040168](https://lccn.loc.gov/05040168) .
\item Shirihai, Hadoram (2008). *A Complete Guide to Antarctic Wildlife* . Princeton, NJ, US: Princeton University Press. [ISBN](/wiki/ISBN\_(identifier)) [978-0-691-13666-0](/wiki/Special:BookSources/978-0-691-13666-0) .
\item Sued-Badillo, Jalil (2003). [*Autochthonous Societies*](https://books.google.com/books?id=zexcW7q-4LgC\&pg=PA65) . General History of the Caribbean. Paris: UNESCO. [ISBN](/wiki/ISBN\_(identifier)) [978-92-3-103832-7](/wiki/Special:BookSources/978-92-3-103832-7) .
\item Thorpe, I. J. (1996). [*The Origins of Agriculture in Europe*](https://books.google.com/books?id=YA-EAgAAQBAJ\&pg=PA68) . New York: Routledge. [ISBN](/wiki/ISBN\_(identifier)) [978-0-415-08009-5](/wiki/Special:BookSources/978-0-415-08009-5) .
\end{itemize}

\section{External links}

**Duck** at Wikipedia's [sister projects](/wiki/Wikipedia:Wikimedia\_sister\_projects)

\begin{itemize}
\item [Definitions](https://en.wiktionary.org/wiki/duck) from Wiktionary
<!-- image -->
\item [Media](https://commons.wikimedia.org/wiki/Anatidae) from Commons
<!-- image -->
\item [Quotations](https://en.wikiquote.org/wiki/Birds) from Wikiquote
<!-- image -->
\item [Recipes](https://en.wikibooks.org/wiki/Cookbook:Duck) from Wikibooks
<!-- image -->
\item [Taxa](https://species.wikimedia.org/wiki/Anatidae) from Wikispecies
<!-- image -->
\item [Data](https://www.wikidata.org/wiki/Q3736439) from Wikidata
<!-- image -->
\end{itemize}

\begin{itemize}
\item [list of books](https://web.archive.org/web/20060613210555/http://seaducks.org/subjects/MIGRATION\%20AND\%20FLIGHT.htm) (useful looking abstracts)
\item [Ducks on postage stamps](http://www.stampsbook.org/subject/Duck.html) [Archived](https://web.archive.org/web/20130513022903/http://www.stampsbook.org/subject/Duck.html) 2013-05-13 at the [Wayback Machine](/wiki/Wayback\_Machine)
\item [*Ducks at a Distance, by Rob Hines*](https://gutenberg.org/ebooks/18884) at [Project Gutenberg](/wiki/Project\_Gutenberg) - A modern illustrated guide to identification of US waterfowl
\end{itemize}

\begin{table}[h]
\begin{tabular}{|l|l|}
\hline
[Authority control databases](/wiki/Help:Authority\_control)  Edit this at Wikidata  <!-- image --> & [Authority control databases](/wiki/Help:Authority\_control)  Edit this at Wikidata  <!-- image --> \\ \hline
National & - [United States](https://id.loc.gov/authorities/sh85039879) - [France](https://catalogue.bnf.fr/ark:/12148/cb119761481) - [BnF data](https://data.bnf.fr/ark:/12148/cb119761481) - [Japan](https://id.ndl.go.jp/auth/ndlna/00564819) - [Latvia](https://kopkatalogs.lv/F?func=direct\&local\_base=lnc10\&doc\_number=000090751\&P\_CON\_LNG=ENG) - [Israel](http://olduli.nli.org.il/F/?func=find-b\&local\_base=NLX10\&find\_code=UID\&request=987007565486205171) \\ \hline
Other & - [IdRef](https://www.idref.fr/027796124) \\ \hline
\end{tabular}
\end{table}

Retrieved from " [https://en.wikipedia.org/w/index.php?title=Duck\&oldid=1246843351](https://en.wikipedia.org/w/index.php?title=Duck\&oldid=1246843351) "

[Categories](/wiki/Help:Category) :

\begin{itemize}
\item [Ducks](/wiki/Category:Ducks)
\item [Game birds](/wiki/Category:Game\_birds)
\item [Bird common names](/wiki/Category:Bird\_common\_names)
\end{itemize}

Hidden categories:

\begin{itemize}
\item [All accuracy disputes](/wiki/Category:All\_accuracy\_disputes)
\item [Accuracy disputes from February 2020](/wiki/Category:Accuracy\_disputes\_from\_February\_2020)
\item [CS1 Finnish-language sources (fi)](/wiki/Category:CS1\_Finnish-language\_sources\_(fi))
\item [CS1 Latvian-language sources (lv)](/wiki/Category:CS1\_Latvian-language\_sources\_(lv))
\item [CS1 Swedish-language sources (sv)](/wiki/Category:CS1\_Swedish-language\_sources\_(sv))
\item [Articles with short description](/wiki/Category:Articles\_with\_short\_description)
\item [Short description is different from Wikidata](/wiki/Category:Short\_description\_is\_different\_from\_Wikidata)
\item [Wikipedia indefinitely move-protected pages](/wiki/Category:Wikipedia\_indefinitely\_move-protected\_pages)
\item [Wikipedia indefinitely semi-protected pages](/wiki/Category:Wikipedia\_indefinitely\_semi-protected\_pages)
\item [Articles with 'species' microformats](/wiki/Category:Articles\_with\_\%27species\%27\_microformats)
\item [Articles containing Old English (ca. 450-1100)-language text](/wiki/Category:Articles\_containing\_Old\_English\_(ca.\_450-1100)-language\_text)
\item [Articles containing Dutch-language text](/wiki/Category:Articles\_containing\_Dutch-language\_text)
\item [Articles containing German-language text](/wiki/Category:Articles\_containing\_German-language\_text)
\item [Articles containing Norwegian-language text](/wiki/Category:Articles\_containing\_Norwegian-language\_text)
\item [Articles containing Lithuanian-language text](/wiki/Category:Articles\_containing\_Lithuanian-language\_text)
\item [Articles containing Ancient Greek (to 1453)-language text](/wiki/Category:Articles\_containing\_Ancient\_Greek\_(to\_1453)-language\_text)
\item [All articles with self-published sources](/wiki/Category:All\_articles\_with\_self-published\_sources)
\item [Articles with self-published sources from February 2020](/wiki/Category:Articles\_with\_self-published\_sources\_from\_February\_2020)
\item [All articles with unsourced statements](/wiki/Category:All\_articles\_with\_unsourced\_statements)
\item [Articles with unsourced statements from January 2022](/wiki/Category:Articles\_with\_unsourced\_statements\_from\_January\_2022)
\item [CS1: long volume value](/wiki/Category:CS1:\_long\_volume\_value)
\item [Pages using Sister project links with wikidata mismatch](/wiki/Category:Pages\_using\_Sister\_project\_links\_with\_wikidata\_mismatch)
\item [Pages using Sister project links with hidden wikidata](/wiki/Category:Pages\_using\_Sister\_project\_links\_with\_hidden\_wikidata)
\item [Webarchive template wayback links](/wiki/Category:Webarchive\_template\_wayback\_links)
\item [Articles with Project Gutenberg links](/wiki/Category:Articles\_with\_Project\_Gutenberg\_links)
\item [Articles containing video clips](/wiki/Category:Articles\_containing\_video\_clips)
\end{itemize}

\begin{itemize}
\item This page was last edited on 21 September 2024, at 12:11 (UTC) .
\item Text is available under the [Creative Commons Attribution-ShareAlike License 4.0](//en.wikipedia.org/wiki/Wikipedia:Text\_of\_the\_Creative\_Commons\_Attribution-ShareAlike\_4.0\_International\_License) ; additional terms may apply. By using this site, you agree to the [Terms of Use](//foundation.wikimedia.org/wiki/Special:MyLanguage/Policy:Terms\_of\_Use) and [Privacy Policy](//foundation.wikimedia.org/wiki/Special:MyLanguage/Policy:Privacy\_policy) . Wikipedia® is a registered trademark of the [Wikimedia Foundation, Inc.](//wikimediafoundation.org/) , a non-profit organization.
\end{itemize}

\begin{itemize}
\item [Privacy policy](https://foundation.wikimedia.org/wiki/Special:MyLanguage/Policy:Privacy\_policy)
\item [About Wikipedia](/wiki/Wikipedia:About)
\item [Disclaimers](/wiki/Wikipedia:General\_disclaimer)
\item [Contact Wikipedia](//en.wikipedia.org/wiki/Wikipedia:Contact\_us)
\item [Code of Conduct](https://foundation.wikimedia.org/wiki/Special:MyLanguage/Policy:Universal\_Code\_of\_Conduct)
\item [Developers](https://developer.wikimedia.org/)
\item [Statistics](https://stats.wikimedia.org/\#/en.wikipedia.org)
\item [Cookie statement](https://foundation.wikimedia.org/wiki/Special:MyLanguage/Policy:Cookie\_statement)
\item [Mobile view](//en.m.wikipedia.org/w/index.php?title=Duck\&mobileaction=toggle\_view\_mobile)
\end{itemize}

\begin{itemize}
\item 
Wikimedia Foundation
<!-- image -->
\item 
Powered by MediaWiki
<!-- image -->
\end{itemize}

\end{document}